\section{Related Work}
\label{sec:related}

\subsection{Document Parsing Models}

Document parsing methods can be roughly divided into multi-stage systems and end-to-end document VLMs. Multi-stage systems, such as classical OCR engines~\cite{paddleocrv2}, the PaddleOCR-VL series~\cite{cui2026paddleocrvl,cui2025paddleocr3}, MinerU~\cite{wang2024mineru} and its decoupled successor MinerU2.5~\cite{niu2025mineru25}, and anchor-prompting parsers such as Dolphin~\cite{feng2025dolphin}, typically perform layout analysis, region cropping, or local recognition before final parsing. They reduce redundant computation through region selection and achieve strong performance on complex documents, but require coordination among multiple modules and are difficult to optimize end to end.

End-to-end document VLMs instead formulate document parsing as unified generation. General-purpose VLM families, such as LLaVA~\cite{liu2023llava}, Qwen-VL~\cite{bai2023qwenvl,wang2024qwen2vl}, InternVL~\cite{chen2024internvl}, MiniCPM-V~\cite{minicpmv2024}, DeepSeek-VL2~\cite{wu2024deepseekvl2}, and Nemotron~\cite{nemotron2026}, acquire strong document understanding through large-scale multimodal pre-training. Building on such backbones, GOT-OCR2.0~\cite{wei2024got}, DeepSeek-OCR~\cite{wei2025deepseekocr} and its successor~\cite{deepseekocr2}, Qianfan-OCR~\cite{dong2026qianfanocr}, dots.ocr~\cite{li2025dotsocr}, HunyuanOCR~\cite{lyu2025hunyuanocr}, FireRed-OCR~\cite{wu2026firered}, and ABot-OCR~\cite{jiang2026abot} further specialize architectures or data for OCR and document scenarios, and recent analyses probe how much such gains rely on the visual versus the language side~\cite{deepseekocr_crutch2026}. To cope with high-resolution pages, mPLUG-DocOwl2~\cite{hu2024docowl2} compresses each page into a compact set of tokens for OCR-free multi-page understanding. However, high-resolution pages still introduce many visual tokens, most of which correspond to blank margins, line spacing, and low-information background regions. Thus, the main bottleneck of end-to-end document VLMs becomes visual token length, motivating a growing body of work on efficient VLMs~\cite{shinde2025surveyefficientvlm,tokencompression_survey2025}.

\subsection{Visual Token Compression}

Visual token compression reduces the visual sequence length fed into LLMs, including fixed compression, token pruning, token merging, and trainable compression. Fixed compression, such as pixel unshuffle, pooling, strided convolution, or fixed projectors, is simple and efficient but content-independent, making it hard to distinguish text-dense regions from blank areas. Learnable projectors partly relax this rigidity: TokenPacker~\cite{li2024tokenpacker} injects high-resolution detail into low-resolution queries for a high compression ratio, and PVC~\cite{yang2024pvc} progressively compresses tokens to unify image and video processing, yet both still apply a compression budget that is fixed at design time rather than adapted to page content.

Token pruning removes low-value tokens using attention scores, feature norms, or learnable importance scores~\cite{rao2021dynamicvit,chen2024fastv,visionzip2025,hu2026illava}. Later methods make the retained set diversity-aware or budget-adaptive: DivPrune~\cite{alvar2025divprune} and IDPruner~\cite{idpruner2026} balance importance and diversity, while others tune the retained count per input or exploit spatial structure~\cite{li2026occamtoken,vispco2026,visiontrim2026,gridprune2025,fcot2025}. For documents specifically, index-preserving pruning~\cite{indexpreserving2025}, reading-inspired schedules~\cite{rtprune2026}, and end-to-end learned selection~\cite{zhu2025visionselector} keep the tokens most useful for recognition. However, discarded information cannot be recovered, which may hurt documents, videos, and fine-grained visual understanding.

Token merging shortens sequences by fusing similar tokens while retaining part of the compressed information. ToMe~\cite{bolya2023tome} uses bipartite matching based on Key-feature similarity, LLaVA-PruMerge~\cite{shang2024llavaprumerge} combines attention-based pruning with similarity-based merging, and recent work merges tokens inside or across encoder layers: TokenCarve~\cite{tan2025tokencarve} preserves information during compression, LaCo~\cite{liu2025laco} performs layer-wise compression within the encoder, and DyMU~\cite{wang2025dymu} dynamically merges and virtually unmerges tokens. In the video domain, where temporal redundancy dominates, a range of schemes~\cite{dycoke2025,shao2025holitom,liu2026infomerge,forestprune2026,onetokenperframe2026,unifiedspatiotemporal2026,smallvlmcompressor2026,statefultoken2026} push compression even higher, but are tuned for temporal rather than spatial density and do not transfer directly to single-page documents.

A complementary line makes the compression policy trainable rather than heuristic. VisionThink~\cite{yang2025visionthink} and explainability-guided task-related compression~\cite{lei2026taskrelated} use reinforcement learning or task signals to decide what to keep, and RL-based token reduction has also been explored for video understanding~\cite{wang2026score,marc2025} and for online token merging in ViTs~\cite{dora2026}. Our RL stage follows this direction and optimizes a lightweight keep/drop policy with GRPO~\cite{shao2024deepseekmath}. Yet existing methods often rely on fixed merge numbers, global thresholds, or schedules designed for temporal redundancy, and do not adapt to the spatial information-density variation within document pages. \ours{} addresses this by using merge-aware SFT to adapt the model to merged visual representations, followed by RL to learn content-adaptive compression boundaries under quality constraints.
